% Based on templates/lecture.tex. Security.pdf, course Week 13.
\begingroup
\renewcommand{\recordid}[1]{}
\renewcommand{\relatedexercise}[2]{\hyperref[exercise:#1]{#2}}
\section{第 24 节课：操作系统保护与安全}
\label{lecture:SC2005-L24}
\lecturemeta
  {SC2005-L24}{2026-11-12（第十三教学周）}
  {Security，OS Protection \& Security}
  {PDF pp.~3--36；Week 13 的 Protection and Security 在本节讲完}
  {Guest Lecture on OS Protection \& Security（仅核对标题）}
  {已完成学习；讲义与例题已核对}

\linespread{1.1}\selectfont
\setlength{\parskip}{0.3em}

\lecturetopic{SC2005-L24-T01}{从失去什么判断安全目标}
手机播放器请求删除照片时，操作系统要检查这个进程对照片的权限。程序的名字和用途无法决定它能做什么。同一批照片也可能面临三种不同的问题（pp.~3、6--10）。

\begin{center}
\small
\begin{tabularx}{\linewidth}{@{}lXX@{}}
\toprule
目标 & 要保护的性质 & 以照片为例 \\
\midrule
机密性（confidentiality） & 只有获准者可以获知内容 & 未获授权的应用读走了一份照片 \\
完整性（integrity） & 内容只能按授权方式更改 & 应用擅自替换了原图 \\
可用性（availability） & 获准者需要时可以使用资源 & 相册服务瘫痪，用户无法取回照片 \\
\bottomrule
\end{tabularx}
\end{center}
同一事件可以影响多个目标。判断“有人偷看，但原文件仍在”时，关键证据是未经授权的读取，受损的是机密性。追查是谁做的属于问责（accountability）；若还要向第三方证明某人确实执行过某项操作，则涉及不可否认性（non-repudiation）。

\lecturetopic{SC2005-L24-T02}{从登录到访问文件：认证、授权与审计}
\textbf{认证（authentication，pp.~12--18）。}输入用户名是在声明身份；系统随后验证凭据，如口令（你知道什么）、令牌（你持有什么）或指纹（你的身体特征）。为了验证口令，系统可以为每个账户生成独立的随机盐值（salt），保存盐值及口令哈希 $F(\mathrm{salt},\mathrm{password})$。登录时取出该账户的盐，用同一函数计算并比较。两个账户即使口令相同，不同盐值通常也会产生不同的结果。盐可以公开，它使一份预计算结果难以在多个账户间复用。

第 14 页区分了哈希函数的两个性质。给定一个摘要，很难找到会产生它的输入，称为\textbf{原像抗性}；很难找到两个不同输入产生同一摘要，称为\textbf{碰撞抗性}。容易猜中的口令即使用了盐仍有风险，因为核对一个猜测只需正向计算。

\textbf{授权（authorization，pp.~19--24）。}系统把请求拆成主体、操作、对象，例如“Carl 要写 \texttt{/usr}”。在访问控制矩阵 $M$ 中，主体占一行，对象占一列。只有请求的操作出现在 $M[\text{主体},\text{对象}]$ 中才放行，空单元格表示没有授予权限。操作系统中负责核对请求并执行决定的机制称为 reference monitor。已通过登录验证的用户，每次请求仍受相应权限约束。

按\textbf{列}整理得到对象的访问控制列表（access control list，ACL），适合回答“谁能改这个目录”；按\textbf{行}整理得到主体的能力列表（capability list），适合回答“这个进程能操作哪些对象”。实际的 capability 是系统保护的对象引用及其权限，持有者不能自行扩大授权。组可以让多个用户共用一项授权。矩阵的读法见\relatedexercise{SC2005-L24-E01}{Alice、Bob 与 Carl 的目录权限}，文件和目录上 \texttt{rwx} 的具体含义见\hyperref[topic:SC2005-L20-T03]{第 20 节}。

\textbf{审计（auditing，p.~25）。}日志记录哪个主体在何时对哪个对象做了什么，以及请求是否成功。它能帮助还原一次文件修改的经过；若日志本身可被当事进程改写，其证明力也会下降。

\lecturetopic{SC2005-L24-T03}{用户程序与内核之间的权限边界}
硬件把普通指令与特权操作分开。用户进程读自己的数据可以直接执行；若要改变受保护的内核状态，就要通过系统调用请求服务（p.~27）。用此前的 xv6 实验来理解，进程调用 \texttt{sbrk} 时陷入内核态，内核检查参数、调整地址空间，再返回用户态。一次系统调用可以始终由同一进程执行。

讲义 pp.~28--31 区分了恶意软件的传播方式和行为。病毒依附宿主传播，蠕虫能够自行传播，木马伪装成正常程序。间谍软件窃取信息，勒索软件限制用户访问自己的数据。Rootkit 隐藏恶意活动并维持特权；一旦内核被控制，依赖它列出进程的检查工具也可能得到不完整的结果。程序能造成多大影响取决于它的权限，例如内核中的越界写可能破坏多个进程共用的数据。

\lecturetopic{SC2005-L24-T04}{谁来验证操作系统}
前面的权限检查都依赖内核。如果启动时加载的内核已被替换，它也能改写检查规则。讲义 pp.~33--35 的办法是从受保护的信任根开始，让已经受信任的一层在移交控制前验证下一层。典型顺序为固件验证引导加载程序，加载程序再验证内核。验证某个镜像的签名时，要用系统信任的公钥，并核对签名是否适用于即将执行的内容。只把普通哈希值和镜像放在同一可改写位置，二者可能被一起替换。见\relatedexercise{SC2005-L24-E02}{签名内核的启动判断}。

\begin{center}
\begin{tikzpicture}[font=\small, >=Latex,
  stage/.style={draw=noteBlue, rounded corners, minimum height=8mm, inner xsep=8pt}]
  \node[stage] (root) {受信任的固件};
  \node[stage, right=19mm of root] (loader) {引导加载程序};
  \node[stage, right=26mm of loader] (kernel) {执行内核};
  \node[below=7mm of kernel] (stop) {拒绝加载};
  \draw[->] (root) -- node[above]{验证通过} (loader);
  \draw[->] (loader) -- node[above]{内核签名有效} (kernel);
  \draw[->] (loader.south) |- node[pos=0.73,above]{验证失败} (stop.west);
\end{tikzpicture}
\end{center}

讲义 p.~36 进一步讨论 enclave，如 Intel SGX。硬件可以为应用的一部分代码和数据提供隔离；远程一方可用远程证明（remote attestation）核对运行环境和代码的度量值，再决定是否交付秘密。操作系统仍能影响进程的调度，因此即使 enclave 内的数据受到保护，服务也可能中断。

\textbf{一分钟回顾。}认证核对身份，授权检查这个身份对指定对象的操作权，审计留下可追查的记录。读访问矩阵时，找某个对象的 ACL 就看列，找某个主体的权限就看行。启动签名把镜像内容和受信任的签名密钥联系起来；镜像中原有的软件缺陷还需要代码审查或测试来发现。
\endgroup
